% ==========================================================================
% INFO 4617 Web Data Science -- Week 7 Friday activity: Triage an issue
% Build: `make` in handouts/, or `latexmk -pdf triage-an-issue.tex` here.
% The figures are annotated from real captures (img/IMAGES.md). A screen
% that needs a signed-in browser is a hand capture: \handcapture prints it
% once img/NAME.png exists, and leaves it out until then.
% ==========================================================================
\makeatletter\def\input@path{{../common/}{handouts/common/}}\makeatother
\documentclass{handout}

\title[INFO 4617 · Week 7 · Friday activities · Triage an issue]{Triage an issue}
\subtitle{Week 7 Friday activity · easy · about 10 minutes an issue}
\author{Brian C. Keegan, Ph.D.}
\institute{Department of Information Science, University of Colorado Boulder}
\date{October 2, 2026}

% A numbered figure: [width]{path from the repo root}{alt text}{caption}
\newcounter{handoutfig}
\newcommand{\activityfigure}[4][\linewidth]{%
  \par\smallskip\refstepcounter{handoutfig}%
  {\centering
   \BeginAccSupp{method=pdfstringdef,Alt={#3}}%
   \includegraphics[width=#1]{#2}%
   \EndAccSupp{}\par}%
  \smallskip{\small\raggedright\textbf{Figure \thehandoutfig.} #4\par}}
% A hand capture: {file in handouts/week-07/img/}{alt text}{caption}
\newcommand{\handcapture}[3]{%
  \IfFileExists{handouts/week-07/img/#1}%
    {\activityfigure{handouts/week-07/img/#1}{#2}{#3}}%
    {\IfFileExists{../../handouts/week-07/img/#1}%
      {\activityfigure{handouts/week-07/img/#1}{#2}{#3}}{}}}

\setlist[enumerate]{leftmargin=*,itemsep=1.5pt,topsep=2pt}
\setlist[itemize]{leftmargin=*,itemsep=1pt,topsep=2pt}

\begin{document}
\maketitle

\begin{frame}{What triage is}
  \begin{columns}[T]
    \begin{column}{0.57\textwidth}
      To triage an issue, you check it against the book as it is today. Then
      you leave one comment that says what should happen next. You change
      nothing in the book.

      Triage makes the other two activities faster. A pull request needs an
      issue that is still a problem, and a reviewer needs to know what the
      issue asks for.

      Your chapter's board in week 7's slides lists the issues to triage. A
      line like \texttt{\#15 → \#59?} asks one question: does pull request
      \#59 fix issue \#15?

      Work in pairs. One partner works in the browser, and the other reads
      this handout aloud.
    \end{column}
    \begin{column}{0.39\textwidth}
      \begin{block}{The three Friday activities}
        \small
        \begin{tabular}{@{}ll@{}}
          \textbf{Triage an issue} & \textbf{easy}\\
          Review a pull request & medium\\
          Issue to pull request & hard
        \end{tabular}
      \end{block}
      \begin{alertblock}{Do not}
        \small
        \begin{itemize}
          \item click \textbf{Close issue};
          \item change the labels or the assignees;
          \item fix the problem here. That is \textit{Issue to pull request}.
        \end{itemize}
        The instructor merges, closes, and resolves conflicts.
      \end{alertblock}
    \end{column}
  \end{columns}
\end{frame}

\begin{frame}{1 · Open the issue}
  \begin{enumerate}
    \item Sign in to GitHub.
    \item Open the issue. Its address is
      \url{https://github.com/cuinfoscience/Web-Data-Science-Book/issues/}
      followed by its number, for example \texttt{.../issues/110}.
    \item Read the answers on the form:
      \begin{itemize}
        \item \textbf{Where were you?} gives the chapter;
        \item \textbf{Which section or heading?} gives the place in the chapter;
        \item the long answers say what went wrong, or what is missing.
      \end{itemize}
    \item Read the comments under the issue. If someone has already triaged
      it, go to the next issue on your board.
  \end{enumerate}
  \handcapture{issue-page.png}{Screenshot of an open issue on the textbook's
    GitHub repository, showing the issue form's answers and the right-hand
    sidebar.}{An open issue: the form's answers, and \textbf{Development} in
    the right-hand sidebar.}
\end{frame}

\begin{frame}{2 · Check it against the book}
  \begin{enumerate}
    \item Open the book: \url{https://cuinfoscience.github.io/Web-Data-Science-Book/}
    \item Go to the chapter that the issue names. Its \textbf{Table of
      contents} ① jumps to the section (Figure~\ref{fig:book}).
    \item Look for the problem. The book changes every week, so the problem
      may already be fixed.
  \end{enumerate}
  \activityfigure[0.64\linewidth]{handouts/week-07/img/book-page_annotated.pdf}%
    {Screenshot of chapter 4 of the book's website. Marker 1 points to the
     Table of contents in the right-hand sidebar. Marker 2 points to Report
     an issue, at the bottom of the sidebar.}%
    {A chapter of the book (here chapter 4). The \textbf{Table of contents} ①
     jumps to a section. \textbf{Report an issue} ② opens the forms that
     readers fill in, so every issue names a chapter and a section.}%
  \label{fig:book}
\end{frame}

\begin{frame}{3 · Look for a pull request, and for a twin}
  \begin{enumerate}
    \item \textbf{A pull request.} On the issue page, look at
      \textbf{Development} in the right-hand sidebar. Then look down the
      issue's timeline: a pull request that names the issue shows as
      \textit{mentioned this}, with its number. Open the pull request, and
      read what it changes.
    \item \textbf{A twin.} Click the \textbf{Issues} tab. Search for a word
      from the title, such as \texttt{traceroute}. Is another open issue
      about the same problem?
  \end{enumerate}
\end{frame}

\begin{frame}{4 · Leave one comment}
  Find the row that matches what you found. Copy its comment, and change the
  numbers and the words.

  \smallskip
  {\small
  \begin{tabular}{@{}>{\raggedright\arraybackslash}p{0.27\linewidth}%
                  >{\raggedright\arraybackslash\ttfamily}p{0.69\linewidth}@{}}
    \toprule
    \textbf{You found} & \normalfont\textbf{Comment}\\
    \midrule
    A pull request fixes it &
      \#59 fixes this when it merges: it adds the output that this issue asks for.\\[3pt]
    A pull request fixes part of it &
      \#59 adds the output. The second half, a note about rate limits, is still open.\\[3pt]
    Another issue reports the same problem &
      Duplicate of \#13\\[3pt]
    The book already fixed it &
      Fixed in the book: "Technical Norms: robots.txt" now shows the output (checked October 2).\\[3pt]
    It is still a problem, with no pull request &
      Still a problem in "Traceroute" (checked October 2). Size: small.\\[3pt]
    You can't tell what it asks for &
      Suggested title: "Gap: Ch. 5 doesn't say what a hop is". Is that what you meant?\\
    \bottomrule
  \end{tabular}\par}

  \medskip
  Use one of three sizes: \textbf{small}, a sentence or a link;
  \textbf{medium}, a paragraph or one code cell; \textbf{large}, a section,
  or more than one file.

  End the comment with your partner's GitHub name:
  \texttt{Triaged with @partner}. Then click \textbf{Comment}.

  \handcapture{issue-comment.png}{Screenshot of the comment box at the bottom
    of an issue, with a triage comment typed in it.}{The comment box at the
    bottom of the issue, with a triage comment typed in. \textbf{Comment}
    posts it. Do not click \textbf{Close issue}.}

  \medskip
  \textbf{Next.} Triage the next issue on your board. When your chapter's
  triage is done, review a pull request (\breakcode{review-a-pull-request.pdf}).
\end{frame}
\end{document}
